\section{What a Digital Governance Framework is}
\label{sec:dgf}

A Digital Governance Framework is the set of decision milestones, or gates, that a significant change
to an enterprise information system passes before it is allowed to proceed. Whatever its domain, every
gate obeys the same mechanics: a gate dossier is prepared by the project team, reviewed in advance by
experts, and presented to a validation body that renders a traced decision. Specialized gates issue
opinions; a general gate consolidates them and arbitrates the commitment of resources.
Table~\ref{tab:gates} lists eight gate types commonly found in large organizations, with examples of the
tasks performed at each.

\begin{table}[!htb]
\centering\footnotesize\renewcommand{\arraystretch}{1.15}
\caption{Eight gate types, what each decides, and examples of its tasks. The list is an example of a
gate set, not a standard: enterprises name, merge, split, and order their gates differently.}
\label{tab:gates}
\begin{tabularx}{\linewidth}{@{}P{2.3cm}P{3.5cm}LP{2.9cm}@{}}
\toprule
Gate & What it decides & Examples of tasks & Validation body\\
\midrule
General (strategic and budget arbitration) & End-of-phase go or no-go; consolidates the opinions of all
other gates. & Check strategic alignment and portfolio priority; update the business case and budget;
release funding by tranche; verify that earlier reservations are lifted; record the decision; at closure,
validate transfer to run and launch benefits tracking. & Executive or investment committee, sponsor, CFO,
CIO, business directors.\\
IT (technology governance and standards) & Fit with technology standards and the enterprise catalog;
operability in run. & Check catalog conformity and process waiver requests, each with a convergence plan; check
that no existing application covers the need; review licences, capacity, hosting, obsolescence, and ITSM integration; validate the support model; prepare the change advisory board. &
CIO, heads of production and development, IT committee.\\
Architecture (urbanization and network constraints) & Coherence of the business, application, data, and
technical architecture with the information system. & Validate the high-level architecture, then the detailed one; review APIs, flows, interfaces, and data ownership; size bandwidth, latency, addressing,
routing, and cloud connectivity; check availability, performance, and reversibility; record structuring
decisions and architecture debt. & Chief architect, architecture committee.\\
Security architecture (cybersecurity and hardening) & Acceptability of cyber risk. & Classify and analyze
risk; review segmentation, filtering rules, encryption, identity and access management, hardening, and
detection; commission intrusion tests; propose acceptance of residual risks. & CISO, security committee,
business risk owner, accreditation authority.\\
Tech readiness (maturity and technical reliability) & Maturity of the technology and aptitude for
production. & Assess technology readiness level and tool stability; review proof-of-concept and pilot results, load tests, and code quality; review the recovery,
resilience, and backup dossier; confirm operational aptitude. & CTO, technical director, CIO,
production release committee.\\
Procurement (purchasing, contracts, suppliers) & Sourcing strategy and supplier commitment. & Decide make
or buy; run the tender; score offers on multiple criteria and total cost; perform third-party due
diligence; prepare supplier contracts and purchasing budgets. & Head of purchasing, commitment committee,
CFO, sponsor.\\
Legal (law and commitments) & Contractual conformity. & Review liability, service-level clauses and
penalties, reversibility, intellectual property, data clauses, labor law, insurance, and signing powers.
& General counsel, corporate secretary, authorized signatories.\\
Compliance (standards, regulation, ethics) & Conformity with regulation and standards. & Check GDPR,
sector standards such as ISO~27001, PCI-DSS, or health-data hosting, NIS2 and DORA, and the AI Act; review
internal control, ethics, accessibility, and archiving. & Chief compliance officer, data protection
officer, compliance committee.\\
\bottomrule
\end{tabularx}
\end{table}

Four levels of involvement recur at every gate: \emph{deciders}, who hold a deliberative vote and answer
for the decision; \emph{dossier owners}, who prepare, present, and defend the dossier; \emph{reviewing
experts}, who analyze it beforehand and issue a formal opinion, favorable, favorable with reservations,
or unfavorable; and \emph{consulted populations}, mobilized according to the nature and criticality of
the project. The validation body then renders one of five decisions: go; go with reservations, with
conditions to be lifted within a set deadline under an action plan; rework, when the dossier is
incomplete or insufficient; suspension, pending an arbitration, a budget, or an external prerequisite;
and no-go. Risk is scored as ordinal impact multiplied by ordinal likelihood, each from 1 to 4; a score
of 12--16 is critical and calls for rejection or competent governance arbitration, and an
above-tolerance acceptance requires a business risk owner and a dated acceptance record
\citep[pp.~38, 98, 100--101]{canale2026}. The person who must fix a defect is therefore distinct from the
person empowered to accept the business risk. Composition is proportional: a simple business SaaS
project does not mobilize the same experts as a critical infrastructure program, job titles vary, and in
mid-size structures one person often holds several roles.

\subsection{Gates are interchangeable: each enterprise composes its own DGF}
\label{sec:notfixed}

A DGF is composed, not prescribed. Each enterprise decides which gates belong to its DGF, on which routes,
in which order, and with which tasks; no gate of Table~\ref{tab:gates} is mandatory, and none is missing
by nature. Gates are interchangeable in a precise sense: they are modular and reconfigurable.
Because every gate has the same form, a dossier prepared, reviewed, and decided, an enterprise can add,
remove, merge, split, or move a gate. It cannot do so freely. A recomposition must preserve the
obligations, the informational prerequisites of each gate, and the authorization relations between
them: a gate that assesses a scope cannot be moved before the gate that defines that scope without
changing its contract, accepting a provisional result, or providing a return. A common form does not
make two contracts identical, and Section~\ref{sec:handoffs} shows how much a gate depends on what
arrives from upstream. Three things therefore vary from one enterprise to
another, and within one enterprise from one project to another.
\emph{The set of gates varies.} Some organizations merge Architecture and Security into one review;
others add gates for data, accessibility, or sustainability. \emph{The order varies.}
Table~\ref{tab:lifecycle} shows one project life cycle in which the same gate is mobilized in several
phases, with a different task each time: Security classifies and analyzes risk at framing, reviews the
security dossier at design, and runs intrusion tests before production; Architecture validates a
high-level architecture, then a detailed one. In agile delivery the gates are kept but lightened: controls
are spread across increments, and only major commitments give rise to a formal passage. \emph{The tasks
vary.} Third-party due diligence sits in Procurement here and in Compliance elsewhere; a licence review
belongs to the IT gate in one enterprise and to Procurement or Legal in another.

\begin{table}[!htb]
\centering\footnotesize\renewcommand{\arraystretch}{1.15}
\caption{One possible positioning of the gates in a project life cycle. The same gate returns in several
phases with different tasks; the general gate closes each phase.}
\label{tab:lifecycle}
\begin{tabularx}{\linewidth}{@{}P{2.7cm}LP{4.0cm}@{}}
\toprule
Project phase & Specialized gates mobilized, with their task in that phase & General gate decision\\
\midrule
1. Opportunity & IT: does a solution already exist in the catalog? Compliance: regulatory sensitivity of
the subject. & Go for study: authorization to frame.\\
2. Framing and feasibility & Architecture: high-level architecture. Security: classification and initial
risk analysis. Tech readiness: maturity of candidate technologies. Procurement: purchasing strategy.
Compliance: applicable regulations. & Investment go: business case and budget envelope validated.\\
3. Design & IT: catalog conformity. Architecture: detailed and network architecture. Security: security
dossier. Procurement: supplier choice. Legal: contracts. Compliance: impact analysis. & Go for build:
construction spending committed.\\
4. Build and acceptance & Tech readiness: tests, quality, operational aptitude. Security: intrusion tests.
IT: operations dossier, change advisory board. Compliance: evidence of conformity. & Go for production:
deployment authorized.\\
5. Deployment and closure & IT: transfer to run. Procurement and Legal: acceptance, warranties,
reversibility. Compliance: recurring controls. & Project review: closure, lessons learned, benefits
tracking.\\
\bottomrule
\end{tabularx}
\end{table}

The gates, routes, and orders used in this paper are therefore an arbitrary selection, taken as examples.
The three routes of Section~\ref{sec:workflows}, Buy, Integrate, and Build, each mobilize a different subset of
the eight gate types in a different order; another enterprise would draw them differently. Nothing in the argument depends on that selection. The
law of gate substitution applies to a gate because of its mechanics, a dossier read and a traced decision returned,
not because of its name, its position, or its task list; it applies equally to gates that do not
appear in this paper. Law~3 asks automation to preserve the composition the enterprise has decided, not
a universal sequence: the enterprise stays free to recompose its DGF, and the automation does not
recompose it to make its own task easier.

A DGF is a \emph{paved road}: a standardized route with known artifacts, decision rules, interfaces,
and escalation paths. Standardization is what makes it governable by humans today, and it is also what
makes it executable by software tomorrow: instead of reconstructing an open-ended organization, an
agent receives a bounded case, a policy version, tools, and an output contract.

The DGF is staffed by specialist populations, not by one department. Table~\ref{tab:populations}
gives the illustrative workforce used throughout the paper. A population can contribute to several
gates, and a gate can require several populations. R1--R9 are resource identifiers, not nine
sequential approvals. Let $P_d$ denote people in department $d$, $K_d$ their useful hours per period,
and $\alpha_{di}$ their allocation to work package $i$. Allocated capacity is
$\sum_d P_d\alpha_{di}K_d$, with $\sum_i\alpha_{di}\le 1$ across all routes and shared work. Thirty architects spending half their useful time on these activities supply fifteen allocated FTE. Approval
and shared-assurance hours are counted once; the sponsors' commitment at the general gate is a
distinct authority rather than an extra finance population.

\begin{table}[!htb]
\centering\footnotesize\renewcommand{\arraystretch}{1.12}
\caption{One illustrative workforce, aligned with the gate types of Table~\ref{tab:gates} and shared
across three generic routes. The contributions are an analytical mapping; the headcounts are synthetic.}
\label{tab:populations}
\begin{tabularx}{\linewidth}{@{}l P{2.7cm} r LLL@{}}
\toprule
ID & Population (gate) & People & Buy (W1) & Integrate (W2) & Build (W3)\\
\midrule
R1 & Procurement & 10 & Sourcing strategy, tender, supplier due diligence. & Inherited supplier contracts and provider evidence. & External components and services, if any.\\
R2 & Legal / Contracts & 10 & Contract, data-processing and service-level clauses. & Inherited obligations, novation, reversibility. & Licences, intellectual property, data-use terms.\\
R3 & Architecture & 30 & Consulted on integration and data flows; no Architecture gate on this route. & Mapping of the inherited system, target integration. & High-level and detailed architecture.\\
R4 & Security Architecture & 20 & Supplier security package, exposure, controls. & Acquisition audit, exposure, risk conditions. & Design assessment and security conditions.\\
R5 & Compliance / Privacy / GRC & 10 & Processing assessment, transfers, sector rules. & Inherited data and regulatory obligations. & Impact assessment and applicable rules.\\
R6 & IT Production and Standards & 25 & Catalog fit, run and support model. & Inventory, rationalization, obsolescence, run takeover. & Catalog conformity, operations dossier, change board.\\
R7 & Tech Readiness & 15 & Pilot results, maturity of the service. & Stability of the inherited platforms. & Tests, load, code quality, production aptitude.\\
R8 & General Gate Support (PMO, portfolio, risk) & 10 & \multicolumn{3}{@{}l@{}}{Gate dossier, reservations register, schedule, and risk cards on every route.}\\
R9 & Finance / Cost Control & 10 & Pricing, subscription cost controls. & Integration-option costs. & Project budget and cost assumptions.\\
\midrule
 & \textbf{Total} & \textbf{140} & \multicolumn{3}{@{}l@{}}{One shared pool, not 140 people for each route.}\\
\bottomrule
\end{tabularx}
\end{table}

No population is the principal actor of a DGF. Buyers and lawyers open a purchase; production and
standards teams open an integration or a build by asking what already exists; architects, the largest
population of the illustrative pool, shape every route; security architects issue one opinion among
several; compliance officers, readiness experts, portfolio, risk, and finance staff assemble the
decision that the general gate takes and that sponsors commit to. The law of gate substitution applies to each of
these roles in the same way, and the examples below deliberately draw on different populations
(Section~\ref{sec:roles}).

\section{Three example routes built from interchangeable gates}
\label{sec:workflows}

Three generic routes serve as examples throughout the paper. \textbf{Buy} acquires a platform or
service from a supplier. \textbf{Integrate} absorbs an information system that already exists, for
example after an acquisition. \textbf{Build} develops a new internal solution. They abstract the three
route types of the author's practitioner framework, a new platform, a merger or acquisition, and a new
project \citep[pp.~17--22]{canale2026}. Figure~\ref{fig:workflows} draws one possible composition for each. They are three
compositions among many: another enterprise could drop Compliance from Integrate, add Tech readiness to
Buy, put Security architecture before IT in Build, or use gates that are not in Table~\ref{tab:gates}
at all.

\begin{figure}[!htb]
\centering
\begin{tikzpicture}[
  box/.style={draw, rounded corners=2pt, minimum width=1.78cm, minimum height=0.85cm,
              align=center, font=\scriptsize},
  arr/.style={-{Latex[length=1.5mm]}, semithick}, node distance=0.25cm]
\node[anchor=west, font=\small\bfseries] at (0,0.78) {W1 --- Buy};
\node[box, fill=yellow!55, anchor=west] (a0) at (0,0) {Business\\need};
\node[box, fill=gray!18, right=of a0] (a1) {Procurement};
\node[box, fill=brown!25, right=of a1] (a2) {Legal};
\node[box, fill=teal!28, right=of a2] (a3) {Compliance};
\node[box, fill=red!30, right=of a3] (a4) {Security};
\node[box, fill=blue!25, right=of a4] (a5) {IT};
\node[box, fill=green!30, right=of a5] (a6) {General};
\foreach \i/\j in {a0/a1,a1/a2,a2/a3,a3/a4,a4/a5,a5/a6} \draw[arr] (\i)--(\j);

\node[anchor=west, font=\small\bfseries] at (0,-0.92) {W2 --- Integrate};
\node[box, fill=yellow!55, anchor=west] (b0) at (0,-1.7) {Existing\\system};
\node[box, fill=blue!25, right=of b0] (b1) {IT};
\node[box, fill=orange!35, right=of b1] (b2) {Architecture};
\node[box, fill=red!30, right=of b2] (b3) {Security\\architecture};
\node[box, fill=brown!25, right=of b3] (b4) {Legal};
\node[box, fill=teal!28, right=of b4] (b5) {Compliance};
\node[box, fill=green!30, right=of b5] (b6) {General};
\foreach \i/\j in {b0/b1,b1/b2,b2/b3,b3/b4,b4/b5,b5/b6} \draw[arr] (\i)--(\j);

\node[anchor=west, font=\small\bfseries] at (0,-2.62) {W3 --- Build};
\node[box, fill=yellow!55, anchor=west] (c0) at (0,-3.4) {New\\project};
\node[box, fill=blue!25, right=of c0] (c1) {IT};
\node[box, fill=orange!35, right=of c1] (c2) {Architecture};
\node[box, fill=red!30, right=of c2] (c3) {Security\\architecture};
\node[box, fill=violet!28, right=of c3] (c4) {Tech\\readiness};
\node[box, fill=green!30, right=of c4] (c5) {General};
\foreach \i/\j in {c0/c1,c1/c2,c2/c3,c3/c4,c4/c5} \draw[arr] (\i)--(\j);
\end{tikzpicture}
\caption{\textbf{Three generic DGF routes, one illustrative configuration.} Yellow: initiating need or
event; grey: Procurement; brown: Legal; teal: Compliance; red: Security, a supplier security review in Buy and Security architecture elsewhere; orange:
Architecture; blue: IT; purple: Tech readiness; green: General gate. Colors identify gate types, not
automation maturity. Each route mobilizes a different subset of the eight gate types of
Table~\ref{tab:gates}, in a different order, and ends at the general gate. Enterprises define their own
gates, orders, and tasks.}
\label{fig:workflows}
\end{figure}

The triggering need, system, or project is not another approval gate. In Buy, the supplier's own
contribution, its proposal, evidence, and commitments, enters through Procurement; it belongs to another
organization and cannot be removed from an end-to-end automation analysis. In this configuration the
three routes contain seventeen gate occurrences after the trigger: six in Buy, six in Integrate, five
in Build. Every gate type appears at least once. The sponsors' commitment is recorded at the general
gate; it is not drawn as a separate box.

The order differs for a reason in each example. Buy examines the supplier's commitments, contract, data
obligations, and security before the run model is fixed; its security gate is a supplier security review,
and no dedicated Architecture gate is mobilized: architects are consulted on integration and data flows
without a gate of their own. Integrate starts from
what exists: IT inventories the inherited system before Architecture and Security assess it, and Legal
and Compliance then qualify the inherited obligations. Build starts by asking whether a solution
already exists in the catalog, and adds a Tech readiness review before production. These orders are
examples (Section~\ref{sec:notfixed}): the enterprise chooses its gates and their order, and what the
law requires is that automation preserve that choice, not these particular routes.

The figure shows nominal forward paths. Reservations, rework, suspension, evidence requests, and dated
conditions remain part of execution, and several gates can review one shared dossier in parallel.
An implementation that always approves, or always refuses to avoid hard cases, fails the function. Full
automation preserves substantive decision quality and the permitted control flow, not just the order of
labels. The same label can also describe different context-sensitive transformations: the security gate in Buy evaluates a supplier's package and commitments; in Integrate, Security
architecture evaluates an inherited system already scoped by IT.

\section{A gate is a contract, not a profession}
\label{sec:contract}

\begin{definition}[Gate-centered work package]
A work package $G_{w,j}$ is the preparation, inquiry, evaluation, decision production, authorized
action, and handoff surrounding position $j$ of route $w$. For versioned state $S$, accessible evidence
$E$, and policies $\Pi$,
\begin{equation}
(S^{+}, D, U) = G_{w,j}(S, E; \Pi, w), \label{eq:gate}
\end{equation}
where $D$ is a disposition and $U$ an action/evidence trace. The procedure may be interactive and
stochastic.
\end{definition}

The input includes more than a PDF: architecture records, observations obtained through tools, a
supplier's reply, costs, identities, outstanding conditions, and a history of negotiations. Requesting
missing evidence is part of the transformation. The output includes more than prose: it can create a
ticket, block a connection, require a redesign, approve a bounded action, or record a refusal. A valid
negative decision is an output; an unsupported declaration that an absent report was verified is not.
Authorization makes an eligible disposition effective; producing it does not confer that authority.

A contract $C_{w,j}$ specifies admissible input histories, acceptable output and action relations,
evidence obligations, quality tolerances, and authority. Several decisions may be defensible. A
replacement need not reproduce a particular person's wording, but it must meet the accepted decision
standard. The reference standard also records disputed judgments instead of treating every historic
human decision as ground truth.

\begin{definition}[Replacement of gate execution]
An implementation replaces a gate when it satisfies $C_{w,j}$ without requiring a person to complete
its ordinary work, exceptions, verification, or per-case authorization. Human participation that
remains necessary is counted as residual execution labor, whether performed internally, by a
contractor, or behind an apparently autonomous service.
\end{definition}

This distinction is established organizational practice, not a newly discovered property of
computation. Stage-Gate separates the work performed in stages from gates that evaluate deliverables, apply
criteria, and decide whether a project proceeds, waits, stops, or returns for further work
\citep[p.~46]{cooper1990}; artifact-centered process modeling represents the records transformed by
business operations \citep{nigam2003}. These antecedents support implementation neutrality. They do
not by themselves prove that every information contract is computable, sufficiently observable,
economically executable, or legally delegable; Sections~\ref{sec:candidate}
and~\ref{sec:premises} give the reasons to expect it for the DGF.

\subsection{Three conditions for complete execution}
\label{sec:conditions}

For a declared case population and contract version, complete execution has three simultaneous
conditions. \emph{Labor:} no necessary human time on any admitted visit, including exceptions and
per-case authorization. \emph{Quality:} independently adjudicated performance satisfies prespecified
noninferiority margins against the human reference and absolute limits on consequential errors, false
refusals, information leakage, and delay. \emph{Mandate:} the decision is institutionally effective,
not merely a recommendation.

A sample containing no human interventions supports, but cannot prove, zero intervention on every
future input. Operational evidence must therefore combine time records, tests, and an architecture
with no required manual step. Assessment by researchers does not itself make operation manual; any
continuing human assessment required to keep the service qualified does count as support. The
distinction is functional, not organizational. A service with autonomous visits but a human
maintenance team has achieved execution closure, not support closure. Initial development is recorded
as a transition investment; continuing updates, incident response, mandated oversight, and outsourced
verification cannot disappear from the account because they have moved to a supplier.

\subsection{One handoff contract and four operating stages}
\label{sec:stages}

A supplier contract can reach Legal as a PDF, Security as an email attachment, and IT production as a
project-manager summary. Each function then extracts the supplier, hosting locations, security
commitments, unresolved clauses, and evidence expiry dates. The interface gain is to perform this
extraction once and attach a provenance trail that the other functions can inspect. The minimal
reusable record contains case and source version, route and gate, claim, evidence locator, verification
status, action owner, risk owner, deadline, and authorization reference:
\par\noindent\begin{minipage}{\linewidth}
\begin{lstlisting}
route: W3
current_gate: TECH_READINESS
recommendation: NOT_READY
claim: load_test_incomplete_and_restoration_not_tested
evidence: {reference: restore-test-report, status: missing}
action_owner_role: production_lead
risk_owner_role: business_service_owner
approval_reference: null
execution_authorized: false
next_action: run_restoration_test_or_request_authorized_exception
\end{lstlisting}
\end{minipage}\par
The record is a schema illustration. Its empty approval reference blocks authorization. A field
marked missing carries information; replacing it with an invented value destroys that information.
This is the difference between a local copilot and a connected process: a copilot saves one reviewer's
reading time, whereas a dependable evidence object changes the work of several downstream functions. A
human can also create that object, which is why Section~\ref{sec:testing} separates the value of
structure from the value of agentic production.

In a \textbf{human pipeline}, specialists prepare and check every dossier. In an \textbf{assisted
pipeline}, agents collect or draft while humans inspect the ordinary path. In an
\textbf{exception-routed pipeline}, eligible work executes automatically, while experts handle
exceptions, sampled review, and retained authorizations. The thesis adds a fourth stage, the
\textbf{closed pipeline}, in which exceptions, verification, and authorization are executed under
mandate as well. The configurations can coexist within a route, and no company is assumed to move
monotonically: incidents can restore manual review. The transition moves the expert from reviewing every ordinary case to governing the system that
reviews them (Section~\ref{sec:implications}), until those tasks are themselves executed by software.
